\documentclass{beamer}
\usepackage{amsmath}
\usepackage{verbatim}
\usepackage{bm}
\usepackage{animate}
\usepackage{amsmath}
\usepackage{tikz}
\usetheme{Madrid}
\usepackage{tikz}
\usepackage{soul}
\usepackage[normalem]{ulem}
\graphicspath{{./figures/}}

\useinnertheme{circles}
%Information to be included in the title page:
\title{Maximum a posteriori (MAP) estimation}
\author{Computational Biology II }
\date{28th April, 2026}

\renewcommand{\insertshortauthor}{Computational Biology II }  
\renewcommand{\insertshorttitle}{Maximum a posteriori (MAP) estimation}
\renewcommand{\insertshortdate}{28th April, 2026} % Define custom text for 3rd block

\usepackage{amsthm}
\newcommand*{\QED}{\hfill\ensuremath{\blacksquare}}





\begin{document}
	\frame{\titlepage}
	
	
	
	\begin{frame}
		\frametitle{Likelihood and log-likelihood}
		\begin{itemize}
			\item Previously, we introduced the likelihood function:
			$$
			L(\bm{\theta}) = p(D \mid \bm{\theta}) = \prod_{i=1}^{n} p(d_i \mid \bm{\theta})
			$$		
			where $\bm{\theta}$ is the set of model parameters, $p(.)$ is a probability density function, and $D = \{ d_1,...,d_n\}$ are i.i.d. observations.
			\pause
			\item Maximum likelihood estimation (MLE):
			$$
			\hat{\bm{\theta}}_{MLE} = \arg\max_{\bm{\theta}} \; L(\bm{\theta})
			$$
			\pause
			\item Maximizing likelihood is equivalent to maximizing log-likelihood, given as 
			$$
			\ell(\bm{\theta}) = \log L(\bm{\theta}) = \sum_{i=1}^{n} \log p(d_i \mid \bm{\theta})
			$$
		\end{itemize}
	
		
	\end{frame}	
	
	\begin{frame}
		\frametitle{Posterior distribution of parameters}
		\begin{itemize}
			\item So far, we focused on $p(D \mid \bm{\theta})$. What about $p(\bm{\theta} \mid D)$?
			\pause
			\item \textbf{Bayes theorem} gives us
			$$
			p(\bm{\theta} \mid D) = \frac{p(D \mid \bm{\theta}) p(\bm{\theta})}{p(D)}
			$$ 	
			where:
			\begin{itemize}
				\item $p(\bm{\theta} \mid D)$ is the \textbf{posterior}.
				\item $p(D \mid \bm{\theta}) = L(\bm{\theta}) $ is the \textbf{likelihood}.
				\item $p(\bm{\theta})$ is the \textbf{prior}.
				\item $p(D)$ is the \textbf{evidence term}.
			\end{itemize}		
		\end{itemize}
	\end{frame}	
	
		
	\begin{frame}
		\frametitle{Maximum a posteriori (MAP)}
		\begin{itemize}
			\item Instead of maximizing likelihood (MLE), we may decide to maximize posterior (MAP):
			$$
			\hat{\bm{\theta}}_{MAP} = \arg\max_{\bm{\theta}} p(\bm{\theta} \mid D)
			$$	
			\pause
			\item Again, it will be useful to maximize its logarithm instead:
			\begin{align*}			
				\log p(\bm{\theta} \mid D)
				&= \log \left(\frac{p(D \mid \bm{\theta}) p(\bm{\theta})}{p(D)}\right) \\
				&= \log	p(D \mid \bm{\theta}) + \log 	p(\bm{\theta})	- \log 	p(D)					
			\end{align*}
			\pause
			\item From that, we get
			$$
			\hat{\bm{\theta}}_{MAP} = \arg\max_{\bm{\theta}}\; \ell(\bm{\theta}) + \log p(\bm{\theta})
			$$
			
		
		\end{itemize}
		
	\end{frame}
		\begin{frame}
			\frametitle{Example: ridge regression}
			\begin{itemize}
				\item Consider a linear model parametrized by $\bm{\beta} = [\beta_0, \beta_1,...,\beta_p]^T$
				\pause
				\item We assume Gaussian prior on the parameters:
				$$
				\beta_i \sim N(0, \sigma^2) \quad \forall i \in \{1,...,p\}
				$$
				where the density function is 
				$$
				p(\beta_i) = \frac{1}{\sqrt{2 \pi \sigma^2}} \exp \left( - \frac{\beta_i^2 }{2\sigma^2} \right)
				$$
				\pause
				\item We have
				$$
				\log p(\beta_i) = \log \left( \frac{1}{\sqrt{2 \pi \sigma^2}} \right) - \frac{\beta_i^2 }{2\sigma^2}
				$$
				\pause
				\item MAP leads to the penalty:
				$$
				\lambda \sum_{i=1}^{p} \beta_i^2, \qquad \lambda = \frac{1}{2\sigma^2}
				$$			
				
			\end{itemize}
\end{frame}
	
	
\end{document}
